\documentclass[11pt]{article}
\usepackage[margin=1in]{geometry}
\usepackage{setspace}
\usepackage{booktabs}
\usepackage{amsmath,amssymb}
\usepackage{graphicx}
\usepackage{svg}
\usepackage{float}
\usepackage[numbers,sort&compress]{natbib}
\usepackage{xurl}
\usepackage[hidelinks]{hyperref}
\graphicspath{{../../../results/figures/}}
\svgpath{{../../../results/figures/}}
\doublespacing
\title{Household and Community Heterogeneity in Childhood Stunting in Ghana}
\author{Gideon Ofosu Kyere\\\small Kwame Nkrumah University of Science and Technology, Kumasi-Ghana\\\small ORCID: 0009-0003-9848-8437\\\small Correspondence: kyereofosu2003@gmail.com}
\date{}
\setlength{\emergencystretch}{2em}
\begin{document}
\maketitle

\begin{abstract}
\noindent\textbf{Objectives:} To estimate residual household and survey-community heterogeneity in childhood stunting in Ghana and assess its sensitivity to measured covariates and modelling choices.

\noindent\textbf{Methods:} We analysed 4,928 children aged 0--59 months in 3,545 households and 616 survey communities from the 2022 Ghana Demographic and Health Survey. Survey-weighted descriptives were independently reproduced. Bayesian logistic models included household and community random intercepts, WASH, and maternal education. After posting the preprint, we audited all posterior parameters and scientific derived quantities in nine separately saved four-chain reproduction fits.

\noindent\textbf{Results:} Survey-weighted stunting prevalence was 17.4\% (95\% CI 15.8--18.9). Retained eight-chain Model 3 summaries gave household and community SDs of 1.23 and 0.39, with household VPC 0.31 and community ICC 0.03. Selected exploratory ORs were 1.50 (95\% posterior interval 1.24--1.82) for boys, 2.91 (2.06--4.21) for ages 24--35 versus 0--5 months, 1.50 (1.15--1.99) for unimproved water, and 0.37 (0.20--0.65) for higher maternal education. The locked household and community SDs had $\hat R$ above 1.01. Local Model 3 and survey weighting passed the strict full-parameter screen; seven reproduction fits failed its $\hat R$ threshold despite adequate ESS and otherwise healthy sampler statistics.

\noindent\textbf{Conclusions:} Stored summaries suggest greater household heterogeneity, but full validation of the model and sensitivity sequence remains incomplete. The audit does not establish that the retained estimates are incorrect; their uncertainty and robustness claims require qualification. The findings are observational and do not identify causal effects.
\end{abstract}

\noindent\textbf{Keywords:} Growth Disorders; Multilevel Analysis; Bayes Theorem; Socioeconomic Factors; Educational Status; Drinking Water; Africa, Western

\section{Introduction}

Ghana's national stunting prevalence has fallen, but childhood stunting remains unevenly distributed. In the 2022 Ghana Demographic and Health Survey (GDHS), about 17\% of children under five were stunted \citep{GSSICF2024}. The national estimate masks clear differences by region, household wealth, and place of residence. It also leaves an important question unanswered: after measured characteristics are taken into account, how much of the remaining variation is shared within households and how much is shared across survey communities?

The distinction matters because children in the same household share many immediate conditions, including food resources, caregivers, sanitation facilities, and economic circumstances. Households in the same survey cluster may also share markets, health services, infrastructure, and environmental exposures. If a model allows for community clustering but not household clustering, some of the similarity among children living together may be assigned to the wrong level.

Previous Ghanaian studies have documented socioeconomic, household, and geographical differences in stunting. Bayesian geostatistical analyses have described spatial variation \citep{AhetoDagne2021,AmoakoJohnson2022}, while multilevel studies have identified household and socioeconomic correlates \citep{AhetoEtAl2015,IddrisuGyabaah2023}. More recent work across 35 sub-Saharan African countries again points to differences by age, sex, maternal education, wealth, and contextual conditions \citep{MoloroEtAl2026}. A 2026 systematic review also found that WASH--stunting associations vary across exposure definitions and settings \citep{WijayantoEtAl2026}, and work using the 2022 GDHS has examined WASH and child nutrition among children born to young mothers \citep{OgumEtAl2026}. Against that background, our focus is narrower: we separate household from survey-community heterogeneity in the eligible under-five sample, then examine how that decomposition behaves as WASH and maternal education are added, using harmonized samples where needed.

We also examine whether the main findings depend on modelling choices that could affect interpretation, including prior scale, treatment of survey weights, missing maternal-education information, categorical versus spline-based age modelling, and binary versus more detailed WASH definitions.

\section{Methods}

\subsection{Data source and analytic sample}

This was a cross-sectional secondary analysis of the Household Member Recode from the 2022 GDHS, accessed under authorization from The DHS Program. Fieldwork was conducted from October 2022 to January 2023. The GDHS used a stratified two-stage cluster sample \citep{GSSICF2024}. Children were eligible for this analysis if they were aged 0--59 months, had slept in the sampled household on the night before the interview (\texttt{hv103 = yes}), and had a valid height-for-age z-score in \texttt{hc70}.

We defined stunting as height-for-age below $-2$ standard deviations from the WHO reference median, corresponding to \texttt{hc70 < -200} in the DHS recode. The resulting core sample contained 4,928 children in 3,545 households and 616 survey communities. For the hierarchical models, community was operationalized as the DHS survey cluster (\texttt{hv001}), and households were uniquely identified by the combination of cluster and household identifiers (\texttt{hv001} and \texttt{hv002}).

\subsection{Covariates}

The core model included child sex, age, household wealth quintile, urban or rural residence, and administrative region. For the primary model, age was grouped as 0--5, 6--11, 12--23, 24--35, 36--47, and 48--59 months.

Drinking-water source (\texttt{hv201}) and sanitation facility (\texttt{hv205}) were grouped into improved and unimproved/no-facility categories using JMP-consistent facility-type definitions \citep{WHOUNICEFJMP2021}. These groupings refer to source or facility type; they do not represent the full JMP service ladders. Maternal education came from \texttt{hc61} and was classified as no education, primary, secondary, or higher. Maternal education was unavailable for 425 children, leaving 4,503 children for the complete-case maternal-education model.

\subsection{Descriptive analysis}

We used the DHS household-member weight (\texttt{hv005}) for descriptive prevalence estimates. Taylor linearization was used for standard errors and 95\% confidence intervals, with \texttt{hv021} as the primary sampling unit and \texttt{hv022} as the sampling stratum. The sum of the rescaled weights in our analytic sample was 4,292.97, which reproduces the official GDHS weighted height-for-age denominator of 4,293 after rounding.

\subsection{Bayesian hierarchical models}

For child $i$ in household $j$ and community $k$, we modelled
\[
Y_{ijk}\sim\text{Bernoulli}(p_{ijk}),
\]
with
\[
\text{logit}(p_{ijk})
=
\alpha+\mathbf{x}_{ijk}^{\top}\boldsymbol{\beta}+u_k+v_{jk},
\]
where $u_k\sim N(0,\tau_c^2)$ is the community random intercept and $v_{jk}\sim N(0,\tau_h^2)$ is the household random intercept. Both random effects were fitted with non-centred parameterizations.

We used $N(0,1.5^2)$ for the intercept, $N(0,1^2)$ for regression coefficients, and Half-Normal$(1)$ priors for the household and community standard deviations. Model 1 contained the core covariates together with household and community random intercepts. Model 2 added WASH. Model 3 added maternal education. We also refitted Model 2 on the 4,503-child maternal-education sample so that the Model 2--Model 3 comparison was not confounded by a change in sample composition.

For interpretation of the random effects, we used a latent logistic residual variance of $\pi^2/3$. The community ICC and household variance partition coefficient were
\[
ICC_c=
\frac{\tau_c^2}{\tau_c^2+\tau_h^2+\pi^2/3},
\qquad
VPC_h=
\frac{\tau_h^2}{\tau_c^2+\tau_h^2+\pi^2/3}.
\]
We also calculated median odds ratios (MORs) as
\[
MOR=\exp\left[\sqrt{2}\Phi^{-1}(0.75)\tau\right].
\]

\subsection{Computation, diagnostics, and sensitivity analyses}

Models were fitted with NUTS in PyMC. For the strengthened final Model 3 diagnostic run, we combined eight independent chains, each with 500 warmup and 500 retained draws, giving 4,000 retained posterior draws. We examined divergences, $\hat R$, effective sample size, BFMI, and tree depth, and used posterior predictive checks to compare replicated and observed prevalence.

We then challenged the main specification in several ways. We refitted Model 3 with tighter and wider priors; fitted a pseudo-posterior using survey weights rescaled to mean one; retained children with missing maternal education as a separate sensitivity category; replaced the age groups with a cubic B-spline; and replaced the binary WASH variables with more detailed water and sanitation categories. Predictive performance of harmonized Models 2 and 3 was compared with PSIS-LOO and Pareto-$k$ diagnostics. As an independent check on residual clustering and fixed-effect robustness, we also fitted logistic generalized…2180 tokens truncated…hains, and no chain saturated the tree depth. Fixed-effect convergence was strong: the largest $\hat R$ among the reported fixed effects was about 1.003, and bulk effective sample sizes exceeded 2,200. Mixing was slower for the two variance parameters. The community and household SDs had $\hat R$ values of about 1.021 and 1.018, with bulk ESS of about 544 and 524. Both $\hat R$ values exceed the commonly used 1.01 guideline, so the retained fit does not meet the strict joint-posterior screen. Its variance estimates and associated substantive interpretations remain provisional.

The full-parameter reproduction audit passed the strict numerical screen for local Model 3 and survey weighting only. Model 1, full and harmonized Model 2, both prior sensitivities, spline age, and detailed WASH failed the $\hat R<1.01$ requirement (maximum $\hat R$ 1.0104--1.0159). All nine had finite diagnostics, bulk and tail ESS above 400, zero divergences, minimum BFMI above 0.3, and no tree-depth saturation. The complete-case local Model 3 had maximum $\hat R=1.009845$, whereas the weighted fit had 1.008136. These passes do not certify the separate locked eight-chain posterior or the entire sensitivity sequence. The supplementary full-posterior audit table reports the results. No saved missing-maternal-education fit was available among the nine audited files.

The posterior predictive distribution reproduced the observed overall prevalence and the prevalence in each of the six age groups. The historical spline-age summaries were similar to the primary summaries, but the available reproduction fit failed the strict joint screen. Its exploratory estimates were: in the spline-age model, the OR was 1.49 (1.12--2.00) for unimproved water, 0.37 (0.20--0.67) for higher maternal education, and 0.35 (0.19--0.61) for richest versus poorest.

The historical detailed-WASH summaries suggested differences between water-source categories, but that reproduction fit also failed the strict joint screen. Surface-water use had an OR of 1.52 (1.10--2.08) relative to improved sources. The estimate for unprotected groundwater was 1.41 (0.90--2.21). For sanitation, neither open defecation [OR 1.23 (0.92--1.64)] nor other unimproved facilities [OR 0.93 (0.68--1.31)] was clearly separated from the improved-facility group.

The survey-weight pseudo-posterior was less precise for water: OR 1.38 (0.90--2.07). This did not reverse the direction of the estimate, but it weakened the strength of the evidence.

The independent GEE checks pointed in the same direction as the hierarchical model. The exchangeable working correlation was 0.154 when clustering by household and 0.024 when clustering by community. With household clustering, the GEE estimate for unimproved water was OR 1.40 (1.12--1.74); replacing age categories with a spline gave OR 1.39 (1.12--1.74). Under the detailed WASH coding, surface water had OR 1.43 (1.12--1.83), whereas unprotected groundwater remained less precise at OR 1.34 (0.93--1.94). These GEE quantities are sensitivity checks rather than substitutes for the Bayesian variance components.

The retained observation-level PSIS-LOO comparison is inconclusive. It cannot establish predictive equivalence or a reliable ranking while joint-posterior diagnostic limitations remain. Model 3 improved ELPD by 0.92, with an SE of 3.38. No Model 3 observation had Pareto $k>0.70$; Model 2 had five, and neither model had $k>1$. Because individual children are left out while household and community effects are estimated from the remaining data, this comparison reflects conditional prediction within the observed hierarchy; it should not be read as performance for entirely new households or communities.

\section{Discussion}

The stored model summaries suggest a larger household than survey-community component after adjustment. This pattern is a hypothesis supported by exploratory estimates, rather than a fully validated result across the entire model and sensitivity sequence.

This pattern is consistent with earlier Ghanaian work showing meaningful household-level variation in childhood malnutrition \citep{AhetoEtAl2015}, and it does not conflict with evidence of geographical heterogeneity from geostatistical studies \citep{AhetoDagne2021,AmoakoJohnson2022}. A DHS cluster captures a broader shared setting, whereas a household captures a more immediate one. Treating the cluster as the only higher level can blur these two sources of dependence.

The median odds ratios make the contrast easier to see: 3.24 at the household level compared with 1.45 at the community level in the final model. These values summarize the spread of the random effects; they are not causal effects of moving a child from one household to another. Even so, the size of the household term suggests that measured DHS covariates leave important shared household conditions unaccounted for. Food security, food allocation, maternal nutrition, caregiving, recurrent infection, and short-term economic shocks are plausible contributors.

The age pattern was clear: stunting was most common around the second and third years of life, and adjusted odds were highest at 24--35 months. The stored spline-age estimates were similar to the categorical-age estimates. Because the available spline reproduction fit failed the full-parameter screen, those similarities are exploratory and do not establish insensitivity to the chosen age cut-points.

The stored Bayesian summaries suggest a wealth gradient, consistent with the verified descriptive pattern; robustness across all model and prior specifications remains provisional. Even after WASH and maternal education were in the model, children in the richest households had markedly lower odds of stunting than those in the poorest. Similar socioeconomic inequalities have been reported repeatedly in Ghana \citep{SeworJayalakshmi2024,OsborneEtAl2025}. The attenuation of the wealth coefficient after maternal education was added is unsurprising given the relationship between schooling and economic position, but these cross-sectional data cannot establish mediation.

The WASH results were less uniform. A binary unimproved-water indicator was associated with higher odds of stunting in the main model, but the weighted pseudo-posterior was more uncertain. The exploratory detailed coding suggested a clearer surface-water signal than for unprotected groundwater; that interpretation remains provisional because the saved reproduction fit failed the joint screen. Sanitation, despite a large crude difference in prevalence, was weak after multivariable adjustment. The evidence for water source is therefore clearer in these data than the evidence for an independent sanitation association. It does not show that changing water source alone would prevent stunting.

Higher maternal education was associated with substantially lower odds of stunting, whereas primary and secondary education were not clearly different from no education once the other covariates were included. Similar inverse associations between maternal schooling and child undernutrition have been reported elsewhere in Ghana and sub-Saharan Africa \citep{AhetoEtAl2015,TakeleEtAl2022,AgyenEtAl2024}. The pattern here should not be read as evidence of a precise educational threshold. Maternal education is also a marker for differences in health literacy, income opportunities, autonomy, healthcare use, and family circumstances that we cannot fully separate in this dataset.

The model comparison also deserves restraint. Maternal education is substantively relevant, but the retained comparison does not establish either a predictive improvement or equivalence. The ELPD difference between harmonized Models 2 and 3 was small relative to its uncertainty. Because the comparison is conditional on the sampled hierarchy, it does not test prediction in entirely new households or communities. The audited harmonized Model 2 reproduction fit failed the joint screen, and the locked Model 3 variance scales exceed 1.01. We therefore retain the historical comparison as provisional and retain maternal education for substantive reasons, without a claim of predictive superiority or equivalence.

\subsection{Strengths and limitations}

The analysis has several practical strengths. The reconstructed descriptive sample agrees with the published GDHS after rounding for both the weighted denominator and national stunting prevalence. The models separate household from community clustering, use a harmonized sample when Models 2 and 3 are compared, and test the main findings under several alternative specifications rather than relying on a single fit.

The study also has important limitations. The GDHS is cross-sectional, so the odds ratios describe associations rather than causal effects. Several plausible determinants of child growth---including detailed diet, household food insecurity, maternal nutritional status, birth size, and repeated infection---were not available in the analysis. Some of that unmeasured information may be reflected in the household random effect.

Our WASH variables describe source and facility types, not the full JMP service ladders. Maternal education was missing for 8.6\% of the core sample. Historical missing-category summaries were similar, but no saved fit for that sensitivity was available for the full-parameter audit, so it remains unaudited. The primary Bayesian likelihood does not fully incorporate the survey design; the weighted pseudo-posterior was therefore treated as a sensitivity analysis.

Finally, household variance is harder to estimate than the fixed effects. Many households contributed only one eligible child, and the two random-effect SDs mixed more slowly than the regression coefficients. The locked eight-chain variance-scale diagnostics exceed 1.01, and seven of nine separately saved reproduction fits fail the strict joint screen. These limitations concern validation of the joint posterior and its derived conclusions, not only exact interval endpoints. The full sensitivity sequence remains provisional. Passing local Model 3 and weighted diagnostics do not establish exact provenance or full validation of the locked posterior. Trace/rank and estimand-specific Monte Carlo error review remains unfinished.

\subsection{Conclusions}

The 2022 GDHS shows that childhood stunting in Ghana remains geographically and socioeconomically uneven, but the variation is not only between places. Stored model summaries suggest greater residual household than survey-community heterogeneity after WASH and maternal education adjustment. That pattern warrants further investigation, but complete joint-posterior and sensitivity validation is needed before treating it as a confirmed variance contrast.

The exploratory contrast suggests that measured DHS covariates may leave shared household conditions unaccounted for; the magnitude and robustness of that contrast require further validation. Future work could examine that residual household component more directly using richer information on food security, diet, maternal health, infection, caregiving, and household shocks.


\section*{Data Availability Statement}
The 2022 Ghana DHS microdata are distributed by The DHS Program and cannot be redistributed by the author. Eligible researchers may request access directly from The DHS Program under its data-use conditions. Reproducibility code, aggregate tables, figures, and non-identifying diagnostic outputs are available at \url{https://github.com/OG-Kyere/Bayesian-Hierarchical-Modelling-of-Childhood-Stunting-in-Ghana}. Restricted row-level DHS data are not included.

\section*{Funding}
The author received no specific funding for this work.

\section*{Conflict of Interest}
The author declares no conflict of interest.

\section*{Author Contributions}
Gideon Ofosu Kyere: Conceptualization; Data curation; Formal analysis; Investigation; Methodology; Software; Validation; Visualization; Writing -- original draft; Writing -- review \& editing; Project administration.

\section*{Word Count}
The main-body and abstract word counts should be measured from the final submission PDF; the structured abstract remains under 300 words.

\section*{Supporting Information}
Additional supporting information contains sample-validation and household-structure summaries, sampler and parameter diagnostics, variance summaries, prior and survey-weight sensitivity analyses, missing-data and age-functional-form sensitivity analyses, detailed WASH and GEE robustness checks, posterior predictive checks, and the PSIS-LOO comparison.

\bibliographystyle{vancouver}
\bibliography{../../../thesis/references}
\end{document}
